\section{Capítulo~\ref{sec:nora}. \nt{NORA}}

A organização do sistema \nt{NORA}, seus dois modos, a ligação com a pupila (não só via \nt{NORA}), o aumento da relação sinal-fundo e o U invertido no comportamento foram medidos diretamente; o ganho multiplicativo $g$ é um modelo; que o ganho comute o modo de escolha e funcione como inverso da temperatura são deduções do capítulo; a vantagem do ajuste do ganho é um cálculo com o modelo do livro; ``o \nt{NORA} é o botão da busca'' e ``o \nt{NORA} é uma interrupção'' são hipóteses.

{\footnotesize
\setlength{\tabcolsep}{3pt}\renewcommand{\arraystretch}{1.15}
\begin{longtable}{>{\raggedright}p{0.5cm}>{\raggedright}p{4.0cm}>{\raggedright}p{5.6cm}>{\raggedright}p{2.3cm}>{\raggedright\arraybackslash}p{1.7cm}}
\toprule
N.º & Proposição & Experimento e o que foi medido & Fonte & Status \\
\midrule\endfirsthead
\toprule
N.º & Proposição & Experimento e o que foi medido & Fonte & Status \\
\midrule\endhead
1 & Os neurônios \nt{NORA} no ser humano são algumas dezenas de milhares, quase todos num único grupo & Cortes seriados do tronco encefálico humano, coloração para tirosina hidroxilase: $53\,900$ células no locus coeruleus e mais $6\,260$ num subnúcleo vizinho. & \cite{Baker1989LC} & medido \\
2 & As saídas se espalham por quase todo o cérebro, mas partes do grupo atendem em parte regiões diferentes & Marcação viral-genética em camundongos: os neurônios \nt{NORA} do locus coeruleus recebem entradas de amplas regiões do cérebro e enviam saídas a amplas regiões do cérebro e da medula espinhal. Em ratos: o grupo que vai à amígdala reforça o aprendizado do medo, e um grupo separado que vai ao córtex pré-frontal, sua extinção. & \cite{Schwarz2015LC, Uematsu2017LC, Poe2020} & medido \\
3 & Modo de fundo: disparos regulares, com taxa de frações a alguns por segundo, que muda com a vigília & Ratos, registro de neurônios do locus coeruleus no ciclo sono-vigília: taxa máxima na vigília, menor no sono de ondas lentas, quase zero no sono REM; as mudanças de taxa antecedem a mudança de estado. & \cite{AstonJonesBloom1981} & medido \\
4 & Modo de rajadas: rajada curta em resposta a um evento importante para a tarefa, mais ou menos no instante da decisão & Macacos, tarefa de encontrar um alvo raro: todos os neurônios do locus coeruleus respondem com rajada só ao alvo, $\sim90\ms$ depois dele e $\sim200\ms$ antes da resposta com a mão; com mau desempenho, as rajadas são mais fracas. Numa tarefa de escolha entre duas opções, a rajada está mais presa à resposta que ao estímulo e ocorre tanto em respostas certas quanto erradas. & \cite{AstonJones1994vigilance, Clayton2004LC} & medido \\
5 & O diâmetro da pupila é um ponto de teste impreciso do \nt{NORA}: ele segue o \nt{NORA}, o \nt{ACET} e outras regiões & Macacos: os disparos do locus coeruleus, até disparos isolados de uma célula, preveem a dilatação da pupila; uma ligação parecida, porém mais tardia, existe nos colículos e no córtex cingulado. Camundongos: as flutuações da pupila acompanham a atividade das saídas noradrenérgicas e colinérgicas no córtex. & \cite{Joshi2016, Reimer2016} & medido \\
6 & A noradrenalina suprime a atividade de fundo mais que a evocada; o ganho multiplicativo~\eqref{eq:gain} é um modelo que reproduz esse efeito & Macacos acordados, córtex auditivo, iontoforese de noradrenalina: a atividade de fundo dos neurônios é mais suprimida que a resposta às vocalizações de coespecíficos, e a relação sinal-ruído cresce. A sigmoide multiplicativa~\eqref{eq:gain} foi tirada de um modelo de rede; a multiplicação literal da inclinação não foi medida. & \cite{Foote1975NE, ServanSchreiber1990} & medido; $g$: modelo \\
7 & O ganho eleva $d'$ só contra o ruído depois do amplificador~\eqref{eq:gainsnr} (proposição~\ref{prop:gainnoise}) & Dedução no capítulo; no modelo de rede, o ganho melhora a detecção do sinal. Não há experimento que separe o ruído antes e depois do amplificador e teste essa diferença. & dedução no capítulo; \cite{ServanSchreiber1990} & teoria \\
8 & A fronteira $g\beta=1$ comuta a rede de escolha de ``pondera'' para ``decidiu''~\eqref{eq:wtagain} (proposição~\ref{prop:gainwta}, fig.~\ref{fig:pitchfork}) & Análise linear de dois grupos que se inibem mutuamente; dedução no capítulo. Não há medições diretas desse limiar no cérebro. & dedução no capítulo & teoria \\
9 & O ganho é o inverso da temperatura da escolha~\eqref{eq:softmax} & Com ruído logístico depois do amplificador, a probabilidade de escolha é um softmax com $T=\sigma/g$; dedução no capítulo. & dedução no capítulo & teoria \\
10 & O \nt{NORA} define quanto buscar: atividade de fundo alta, $g$ efetivo pequeno (busca); rajada no instante da decisão, $g$ grande (decisão) & Humanos: antes de uma escolha ao acaso, a pupila basal é mais larga que antes de escolher a melhor opção conhecida; pessoas com pupila mais larga tendem a buscar mais. Ratos: a passagem para o modo de escolha aleatória é controlada pela entrada do \nt{NORA} no córtex cingulado anterior. Que a rajada eleve o $g$ efetivo não foi medido diretamente. & \cite{JepmaNieuwenhuis2011, Tervo2014stochastic, AstonJones2005} & hipótese \\
11 & A recompensa depende de $g$ como um U invertido; o melhor $g$ é menor num mundo que muda depressa (proposição~\ref{prop:explore}, fig.~\ref{fig:invertedU}) & \texttt{models/nora\_explore.py}, $60$ execuções de $4000$ tentativas. Mudança a cada $1000$ tentativas: melhor $g=16$, fração $0.976$; a cada $20$: $g=8$, fração $0.886$; com $g=0.5$, $0.77$. O recálculo coincidiu com o capítulo. & \texttt{models/\allowbreak nora\_\allowbreak explore.py} & modelo \\
12 & U invertido no comportamento: melhor com atividade de fundo moderada e rajadas nítidas & Macacos, mesma tarefa da linha~4: nos períodos de bom desempenho, a taxa de fundo é moderada e as rajadas ao alvo são nítidas; com taxa de fundo alta não há rajadas e há mais respostas falsas. As mudanças da atividade de fundo e evocada acompanham as oscilações do desempenho. & \cite{Usher1999LC, AstonJones2005} & medido \\
13 & O \nt{NORA} informa a incerteza inesperada; o \nt{ACET}, a esperada & Teoria de inferência bayesiana com dois tipos de incerteza; nos trabalhos citados não há experimento direto que separe esses sinais por neurotransmissor. & \cite{YuDayan2005} & hipótese \\
14 & Depois de uma mudança brusca, as pessoas aprendem mais depressa, e a pupila se dilata & Humanos, tarefa de predição com mudanças súbitas: variações curtas da pupila acompanham a probabilidade estimada de mudança e, junto com a pupila basal, preveem o quanto os novos dados mudam a estimativa (a taxa de aprendizado); uma variação da pupila não ligada à luz nem à tarefa muda essa taxa. Que aqui a pupila mostre justamente o \nt{NORA} é uma inferência a partir dos dados indiretos da linha~5. & \cite{Nassar2012} & medido \\
15 & A rajada funciona como interrupção: a rede zera seu estado & Não há experimento direto em que uma rajada do \nt{NORA} zere o estado da rede; só foram medidas rajadas em eventos importantes (linha~4). & \cite{BouretSara2005} & hipótese \\
16 & O ajuste de $g$ ou da taxa de aprendizado pela surpresa quase não supera as melhores configurações constantes & \texttt{models/nora\_explore.py}, mundo com épocas de $1000$ tentativas (mudança a cada $1000$ e a cada $20$). Melhor constante $g=8$: $0.925$; ajuste de $g$: até $0.927$; ajuste da taxa de aprendizado: até $0.926$; taxa constante $0.4$: $0.928$. O recálculo coincidiu com o capítulo. & \texttt{models/\allowbreak nora\_\allowbreak explore.py} & modelo \\
\bottomrule
\end{longtable}}
